\documentclass[10pt,a4paper]{amsart}
\usepackage[margin=19mm]{geometry}
\usepackage{amsmath,amssymb,mathtools}
\usepackage[scr=rsfs,cal=euler]{mathalfa}
\usepackage{xfrac}
\usepackage[unicode,hidelinks,bookmarks=false]{hyperref}
\newcommand{\ie}{i.e.\ }
\newcommand{\cf}{cf.\ }
\newcommand{\resp}{respectively\ }
\newcommand{\Subsection}{\subsection}
\providecommand{\nobreakdash}{\nobreak\hbox{-}}
\setlength{\emergencystretch}{2em}
\multlinegap0pt
\usepackage{mdframed}
\usepackage{mdframed}
\usepackage{mdframed}
\usepackage{mdframed}
\usepackage[shortlabels]{enumitem}
\usepackage{amssymb}
\usepackage{xcolor}
\usepackage{mdframed}
\usepackage{float}
\newtheorem{lemma}{Lemma}
\newtheorem{claim}{Claim}
\newtheorem{claimproof}{Claimproof}
\DeclareMathOperator{\Hom}{Hom}
\title{A Concise Proof of the $L_0$ Dichotomy}
\author{Tonatiuh Matos-Wiederhold}
\date{Source: arXiv:2604.02589v1 (CC BY 4.0)}
\begin{document}
\maketitle

\noindent\emph{Source and licence.} A Concise Proof of the $L_0$ Dichotomy. Version arXiv:2604.02589v1 (CC BY 4.0), distributed by arXiv under the Creative Commons Attribution 4.0 International licence, http://creativecommons.org/licenses/by/4.0/, as recorded in the arXiv licence metadata for this exact version and shown on the abstract page https://arxiv.org/abs/2604.02589v1. Full licence notice: \texttt{licenses/arxiv-2604.02589-CC-BY-4.0.txt}.\par\smallskip

\noindent\textit{Editorial scope.} Counted unit: the article's lemma lemma:Ln\_basics (source0.tex lines 407--415). The article's complete original proof of it is delivered with it (source0.tex lines 417--480, one \texttt{proof} environment of 64 lines). No mathematics is written, repaired or re-derived: the statement and the proof are the article's own lines, in the article's own notation, and no retained line is substituted or re-flowed.  Retained uncounted as the source-local context the proof needs: lines 152--154.  Nothing was omitted from any retained span, so the package carries no omission record.  No retained line contains a citation, so the wrapper carries no bibliography and no bibliography entry.  No retained line contains a figure environment, an image-inclusion command or drawing code, so no image is delivered and the package declares no asset dependency.  Editorial formatting only, in addition: an AMS \texttt{amsart} preamble that restates the article's own declarations and macros used by the retained lines, namely the environments \texttt{lemma}, \texttt{claim}, \texttt{claimproof} on separate counters, together with the standard packages those lines need.  The retained environments print in this document as Lemma 1, Claim 1, Claimproof 1 in order of appearance, whereas the article prints the same items with the numbering of its own full document.  The complete native source of the article as distributed in the arXiv e-print for this exact version is bundled unchanged as \texttt{sources/197749/source0.tex}.
\begin{lemma}[First Reflection Theorem]\label{lemma:first-reflection}
    If $\Phi$ is as in the preceding definition, then for any $\Sigma^1_1$ set $A\in\Phi$, there is a Borel set $B\subseteq X$ such that $A\subseteq B$ and $B\in\Phi$.
\end{lemma}

\begin{lemma}\label{lemma:Ln_basics}~ 
    \begin{enumerate}[label=(\alph*)]
        \item Suppose that $n>0$, that $c$ only takes odd values, that $k$ is a natural and $t\in2^{<n}$ is such that $(p_k)^\smallfrown t^\smallfrown(i)$, for $i<2$, are vertices of $L^c_n$ (see (iii) above). Then, they are an odd distance apart in $L^c_n$.

        \item For any $n<\omega$, if $\mathcal L\subseteq\Hom(L^c_n,G)$ is a large Borel set, then for any $N<\omega$ there exists $d\in\omega^\omega$ such that $d\restriction n=c\restriction n$, $d(n)\geq N$ is odd, and $\mathcal L^{+d}$ is a nonempty subset of $\Hom(L^{d}_{n+1},G)$.

        \item For any $n<\omega$, if $\mathcal L\subseteq\Hom(L^c_n,G)$ is a large Borel set, then there exists $d\in\omega^\omega$ such that $d\restriction n=c\restriction n$ and $\mathcal L^{+d}$ is a large subset of $\Hom(L^{d}_{n+1},G)$.
    \end{enumerate}
\end{lemma}

\begin{proof}
    For (a), observe that both vertices come from the same vertex $(p_k)^\smallfrown t$ in $L^c_{|t|+1}$, and so their distance in $L^c_n$ is twice the distance from $(p_k)^\smallfrown t$ to $e_1^{|t|}$ in $L^c_{|t|+1}$ plus the length of the joining path at that stage, which is $c(|t|)+2$.
    Since $c(|t|)$ is odd, $c(|t|)+2$ is odd, and hence their total distance is even plus odd, which is odd.

    For (b), since $\mathcal L$ is large, it is not tiny.
    Hence, for every vertex $u\in V(L^c_n)$ and every natural $k$, $\Phi(\mathcal L(u),k)$ fails.
    In particular, fix the gluing vertex $u_0:=e_1^n\in V(L^c_n)$ (recall that $e_1^0=p_0$).
    Then for every $k$ there exist $\varphi_0,\varphi_1\in\mathcal L$ and an odd walk $(v_0,v_1,\dots,v_m)$ in $G$ with $v_0=\varphi_0(u_0)$, $v_m=\varphi_1(u_0)$, and $m>2k-1$.
    Choosing $k$ large enough, I can ensure that $m\geq N+2$ and $m$ is odd.
    Since $m$ is odd, $d(n):=m-2$ is also odd and satisfies $d(n)\geq N$.
    For each $1\leq j\leq m$, since $(v_{j-1},v_j)\in E(G)=\mathrm{im}(\pi)$, I pick $e_j\in E$ with $\pi(e_j)=(v_{j-1},v_j)$.
    Setting $d(x)=c(x)$ for all $x\neq n$ and noting $L^c_n=L^d_n$ (since $d\restriction n=c\restriction n$), I define $\varphi\in\Hom(L_{n+1}^d,G)$ by $\varphi^0=\varphi_0$, $\varphi^1=\varphi_1$, $\varphi(p_i)=v_{i+1}$ for $0\leq i\leq m-2=d(n)$, and the edge assignments $\varphi(u_0^\frown (0),p_0)=e_1$, $\varphi(p_j,p_{j+1})=e_{j+1}$ for $0\leq j<m-2$, and $\varphi(p_{m-2},u_0^\frown (1))=e_m$.
    Then $\varphi\in\mathcal L^{+d}$.

    I prove (c) by contradiction.
    Recall that constructing $L^c_{n+1}$ from $L^c_n$ only requires knowing the value of $c(n)$.
    Suppose that for all $d$ that agree with $c$ except possibly in the $n$-th value, $\mathcal L^{+d}=\bigcup_m\mathcal F^d_m$ where each $\mathcal F^d_m$ is a tiny Borel subset of $\Hom(L^d_{n+1},G)$.
    For each $m$ and $d$, let $w^d_m\in V(L^d_{n+1})$ and $k^d_m<\omega$ be such that $\Phi(\mathcal F^d_m(w^d_m),k^d_m)$.
    
    Since $w^d_m\in V(L^d_{n+1})$, it is either a \emph{non-path vertex} of the form $u^\smallfrown(i)$ for some $u\in V(L^c_n)$ and $i<2$, or a \emph{path vertex} $p_j$ for some $0\leq j\leq d(n)$ in the joining path.
    I claim that in either case, there exists a vertex $u^d_m\in V(L^c_n)$ and $\ell^d_m<\omega$ such that $\Phi(\mathcal F^d_m({u^d_m}^\smallfrown(i^d_m)),\ell^d_m)$ for some $i^d_m<2$ where $u^d_m{}^\smallfrown(i^d_m)$ is a non-path vertex of $L^d_{n+1}$.

    Indeed, if $w^d_m=u^\smallfrown(i)$ is already non-path, take $u^d_m=u$, $i^d_m=i$, $\ell^d_m=k^d_m$.
    If $w^d_m=p_j$ is a path vertex of $L^d_{n+1}$, then the joining path connects $e_1^n{}^\smallfrown(0)$ to $e_1^n{}^\smallfrown(1)$ via $p_0,\dots,p_{d(n)}$.
    Since $p_j$ and $e_1^n{}^\smallfrown(0)$ are at distance $j+1$ in $L^d_{n+1}$, every homomorphism $\varphi\in\mathcal F^d_m$ maps them to vertices that are at walk-distance $j+1$ in $G$.
    Given any odd walk $(a_0,\dots,a_r)$ in $G$ with $a_0,a_r\in\mathcal F^d_m(e_1^n{}^\smallfrown(0))$, witnessed by $\varphi,\varphi'\in\mathcal F^d_m$, the walk $(\varphi(p_j),\dots,\varphi(e_1^n{}^\smallfrown(0))=a_0,\dots,a_r=\varphi'(e_1^n{}^\smallfrown(0)),\dots,\varphi'(p_j))$ has length $r+2(j+1)$, which has the same parity as $r$ (hence is also odd), and has endpoints in $\mathcal F^d_m(p_j)$.
    By $\Phi(\mathcal F^d_m(p_j),k^d_m)$, $r+2(j+1)\leq 2k^d_m-1$, so $r\leq 2(k^d_m-j-1)-1$.
    It follows that $\Phi(\mathcal F^d_m(e_1^n{}^\smallfrown(0)),\max\{0,k^d_m-j-1\})$.
    Thus, set $u^d_m=e_1^n$, $i^d_m=0$, $\ell^d_m=\max\{0,k^d_m-j-1\}$.
    
    \begin{claim}
        For each $m$ and $d$, there is a Borel set $B^d_m\supseteq\mathcal F^d_m({u^d_m}^\smallfrown(i^d_m))$ such that $\Phi(B^d_m,\ell^d_m)$.
    \end{claim}

    \begin{claimproof}
        I use the First Reflection Theorem (Lemma~\ref{lemma:first-reflection}).
        Fix $\ell=\ell^d_m$.
        It suffices to verify that $\Phi_\ell:=\{A\subseteq X:\Phi(A,\ell)\}$ is $\Pi^1_1$ on $\Sigma^1_1$.
        Let $Y$ be Polish and $A\subseteq Y\times X$ be $\Sigma^1_1$.
        Fix a Polish space $N$ and a continuous surjection $g\colon N\to A$.
        For each odd $r\geq 2\ell+1$, define
        \begin{align*}
            C_r:=\{(y,x_0,\dots,x_r,n_0,n_1,\hat e_1,\dots,\hat e_r)\in Y\times X^{r+1}\times N^2\times E^r:{}\\
            g(n_0)=(y,x_0)\land g(n_1)=(y,x_r)\land\textstyle\bigwedge_{j<r}\pi(\hat e_{j+1})=(x_j,x_{j+1})\}.
        \end{align*}
        Then $C_r$ is closed (by continuity of $g$ and $\pi$), and $D_r:=\mathrm{proj}_Y(C_r)$ is $\Sigma^1_1$.
        Hence $D:=\bigcup_{r\text{ odd},\,r\geq 2\ell+1}D_r$ is $\Sigma^1_1$, and
        $A_{\Phi_\ell}=\{y\in Y:A_y\in\Phi_\ell\}=Y\setminus D$
        is $\Pi^1_1$.
        Since $\mathcal F^d_m({u^d_m}^\smallfrown(i^d_m))$ is analytic and satisfies $\Phi(\cdot,\ell)$, the First Reflection Theorem yields the desired Borel superset $B^d_m$.
    \end{claimproof}

    Take $$\mathcal H^d_m=\{\varphi\in\mathcal L:\varphi(u^d_m)\in B^d_m\},$$ which is a Borel subset of $\Hom(L^c_n,G)$.
    I claim that $\mathcal H^d_m$ is tiny.
    Indeed, $\mathcal H^d_m(u^d_m)=\{\varphi(u^d_m):\varphi\in\mathcal L\land\varphi(u^d_m)\in B^d_m\}\subseteq B^d_m$, so $\Phi(\mathcal H^d_m(u^d_m),\ell^d_m)$.
    
    Now, the set of pairs $(u^d_m,i^d_m)$ ranges over the finite set $V(L^c_n)\times2$ (which does not depend on $d$), and $(m,d)$ range over countably many values.
    Thus $\bigcup_{m,d}\mathcal H^d_m$ is a countable union of tiny sets, hence small.
    The set $\mathcal L_-:=\mathcal L\setminus\bigcup_{m,d}\mathcal H^d_m$ is therefore large (and Borel).
    By part (b), there exist $\varphi$ and $d$ such that $\varphi\in\mathcal L_-^{+d}\subseteq\mathcal L^{+d}$.
    Then $\varphi\in\mathcal F^d_m$ for some $m$.
    It follows that $\varphi^{i^d_m}(u^d_m)=\varphi({u^d_m}^\smallfrown(i^d_m))\in\mathcal F^d_m({u^d_m}^\smallfrown(i^d_m))\subseteq B^d_m$, hence $\varphi^{i^d_m}\in\mathcal H^d_m$.
    But $\varphi\in\mathcal L_-^{+d}$ implies $\varphi^{i^d_m}\in\mathcal L_-$, contradicting $\mathcal L_-\cap\mathcal H^d_m=\emptyset$.
\end{proof}

\end{document}
